\documentclass[10pt,a4paper]{amsart}
\usepackage[margin=19mm]{geometry}
\usepackage{amsmath,amssymb,mathtools}
\usepackage[scr=rsfs,cal=euler]{mathalfa}
\usepackage{xfrac}
\usepackage[unicode,hidelinks,bookmarks=false]{hyperref}
\newcommand{\ie}{i.e.\ }
\newcommand{\cf}{cf.\ }
\newcommand{\resp}{respectively\ }
\newcommand{\Subsection}{\subsection}
\providecommand{\nobreakdash}{\nobreak\hbox{-}}
\setlength{\emergencystretch}{2em}
\multlinegap0pt
\usepackage{bm}
\usepackage{bbm}
\usepackage{dsfont}
\usepackage{xspace}
\usepackage{cancel}
\usepackage{mathtools}
\usepackage{amsthm}
\makeatletter
\@ifundefined{theorem}{\newtheorem{theorem}{Theorem}}{}
\makeatother
\newcommand{\D}{\mathrm{D}}
\newcommand{\Mcal}{\mathcal{M}}
\newcommand{\pw}{\mathrm{pw}}
\renewcommand{\div}{\mathrm{div}}
\renewcommand{\d}[1]{\,\mathrm{d}#1}
\title[Local discontinuous Galerkin FEM for convex minimization]{Local discontinuous Galerkin FEM for convex minimization}
\author{Carsten Carstensen, Ngoc Tien Tran}
\date{Source: arXiv:2604.08246v2 (CC BY 4.0)}
\begin{document}
\maketitle

\noindent\emph{Source and licence.} Local discontinuous Galerkin FEM for convex minimization. Version arXiv:2604.08246v2 (CC BY 4.0), distributed under the Creative Commons Attribution 4.0 International licence, as recorded in the arXiv licence metadata for this exact version and shown on the abstract page https://arxiv.org/abs/2604.08246v2. Full rights notice: licenses/arxiv-2604.08246-CC-BY-4.0.txt.\par\smallskip

\noindent\textit{Editorial scope.} Counted unit: the Theorem whose source label is \texttt{thm:weak-duality} in the article (source0.tex lines 469--474, source label \texttt{thm:weak-duality}), delivered together with the article's own complete original proof of it (source0.tex lines 476--521, one \texttt{proof} environment of 46 lines). No mathematics is written, repaired or re-derived: statement and proof are the article's own text line for line. Required context retained uncounted: def:ldg-dicrete-gradient (lines 365--368); def:stabilization-ldg (lines 389--396); def:div-rec (lines 439--441). Every definition, display and label of those lines is retained unchanged. Omitted from this package and not redistributed: 1--364, 369--388, 397--438, 442--468, 475--475, 522--1228. Nothing inside a retained excerpt is omitted or altered: every retained body is delivered exactly as the article's lines are written, so no omission receipt of any kind accompanies this package. Citations: the retained excerpts carry no citation command at all, so no bibliography entry of the article is transcribed and no reference is redistributed. Reference adaptation: none was needed and none was made; every reference inside the delivered unit resolves inside it, so no unresolved reference or citation is produced. Printed slips kept literal: no printed word, symbol, hypothesis or number of the counted statement or of its proof was altered. Layout and class: the article's own class is \texttt{amsart} with packages including inputenc, babel, amsfonts, amsmath, amssymb, amsthm, mathtools, esint, tikz, enumitem, hyperref, cleveref, thmtools, xcolor; the wrapper is the lane's pinned \texttt{amsart} editorial wrapper at 10pt on A4, which restates the theorem-like environments the retained text needs and adds the article's own macro definitions that the retained text needs (\texttt{\textbackslash D}, \texttt{\textbackslash Mcal}, \texttt{\textbackslash d}, \texttt{\textbackslash div}, \texttt{\textbackslash pw}), with extra wrapper packages (bm, bbm, dsfont, xspace, cancel, mathtools). The delivered numbering is the unit's own. Assets: no retained span contains a figure environment, a graphics-inclusion command, a PDF-page inclusion, a table or a TikZ picture; no image file is copied into this package, the package declares no asset dependency and no image file is redistributed with it. Licence: version arXiv:2604.08246v2 of this article is distributed under the Creative Commons Attribution 4.0 International licence, as recorded in the arXiv licence metadata for this exact version and shown on the abstract page https://arxiv.org/abs/2604.08246v2. A full rights notice is delivered at licenses/arxiv-2604.08246-CC-BY-4.0.txt. Characters: every retained excerpt is pure ASCII, so the delivered engine, XeLaTeX with its default Latin Modern text font, reports no missing character.
\begin{align}\label{def:ldg-dicrete-gradient}
	\int_\Omega \nabla_h v_h \cdot \Phi \d{x}
	= \int_\Omega \nabla_\pw v_h \cdot \Phi \d{x} - \sum_{S \in \mathcal{F}\setminus\mathcal{F}_\mathrm{N}} \int_S [v_h]_S \{\Phi\}_S\cdot \nu_S \d{s}.
\end{align}

\begin{align}\label{def:discrete_energy}
	E_h(v_h) &\coloneqq \int_\Omega (W(\nabla_h v_h) 
    + \psi_h(\bullet,v_h)) \d{x} + s_h(v_h)/r,\\
    s_h(v_h; w_h) &\coloneqq 
    \sum_{S \in \mathcal{F} \setminus \mathcal{F}_\mathrm{N}} 
    h_S^{-s}\int_S |[v_h]_S|^{r-2} [v_h]_S\,[w_h]_S \d{s}
\label{def:stabilization-ldg}
\end{align}

\begin{align}\label{def:div-rec}
	\int_\Omega \div_h \tau\, \phi \d{x} = -\int_\Omega \tau_\Mcal \cdot \nabla_\pw \phi \d{x} + \sum_{S \in \mathcal{F} \setminus \mathcal{F}_\mathrm{N}} \int_S \tau_S [\phi]_S \d{s} 
\end{align}

\begin{theorem}[duality of LDG]\label{thm:weak-duality}
	It holds $\sup E_h^*(Y) \leq \min E_h(V_h)$; (B1) and (B5) imply 
	%\begin{align}%\label{eq:strong-duality-ldg}
		$\max E_h^*(Y) = \min E_h(V_h)$.
	%\end{align}
\end{theorem}

\begin{proof}
	Given $\tau = (\tau_\Mcal, \tau_\mathcal{F}) \in Y$ and $v_h \in V_h$,
    \eqref{def:ldg-dicrete-gradient} implies
    \begin{align*}
        &\int_\Omega \Pi_\Mcal^{k-1} \tau_\Mcal \cdot \nabla_h v_h \d{x}\\
        &\qquad= \int_\Omega \nabla_\pw v_h \cdot \Pi_\Mcal^{k-1} \tau_\Mcal \d{x} - \sum_{S \in \mathcal{F}\setminus\mathcal{F}_\mathrm{N}} \int_S [v_h]_S \{\Pi_\Mcal^{k-1} \tau_\Mcal\}_S \cdot \nu_S \d{s}.
    \end{align*}
    Since $\Pi_\Mcal^{k-1}$ can be omitted in the integral $(\nabla_\pw v_h, \Pi_\Mcal^{k-1} \tau_\Mcal)_{L^2(\Omega)}$,
    this and the definition of the divergence reconstruction in \eqref{def:div-rec} provide
    \begin{align}\label{eq:diby-hho}
        \int_\Omega \tau_\Mcal \cdot \nabla_h v_h \d{x} &= \int_\Omega \Pi_\Mcal^{k-1} \tau_\Mcal \cdot \nabla_h v_h \d{x} = -\int_\Omega \div_h \tau \,v_h \d{x}\nonumber\\
		&\qquad+ \sum_{S \in \mathcal{F} \setminus \mathcal{F}_\mathrm{N}} \int_S [v_h]_S (\tau_S - \{\Pi_\Mcal^{k-1} \tau_\Mcal\}_S \cdot \nu_S) \d{s}.
    \end{align}
	From this, $\tau_\Mcal \cdot \nabla_h v_h \leq W(\nabla_h v_h) + W^*(\tau_\Mcal)$ and $\div_h \tau_h \, v_h \leq \psi_h(\bullet,v_h) + \psi_h^*(\bullet,\div_h \tau)$ a.e.~in $\Omega$ as well as the H\"older inequality, we deduce $\sup E_h^*(Y) \leq \min E_h(V_h)$.
    To establish equality,
	let $u_h \in \arg\min E_h(V_h)$ be a minimizer of $E_h$ in $V_h$. We define the dual variable $y = (\sigma_\Mcal, \sigma_\mathcal{F}) \in Y$ with
	\begin{align}\label{def:dual-variable}
			\sigma_\Mcal \coloneqq \D W(\nabla_h u_h) \quad\text{and}\quad
			\sigma_\mathcal{F}|_S \coloneqq \{\Pi_{\Mcal}^{k-1} \sigma_\Mcal\}_S \cdot \nu_S - h_S^{-s}|[u_h]_S|^{r-2} [u_h]_S
	\end{align}
	for any $S \in \mathcal{F}\setminus\mathcal{F}_\mathrm{N}$.
	The Euler-Lagrange equations read
	\begin{align}\label{eq:dELE}
		0 = \int_\Omega (\sigma_\Mcal \cdot \nabla_h v_h \d{x} + \partial_u \psi_h(\bullet, u_h) v_h) \d{x} + s_h(u_h;v_h)
	\end{align}
	for any $v_h \in V_h$.
	This and
	the discrete
	integration by parts formula \eqref{eq:diby-hho} imply
	\begin{align}\label{eq:proof-mldg-equivalence}
		0 = &\int_\Omega (\partial_u \psi_h(\bullet, u_h) - \div_h y) v_h \d{x}\nonumber\\
		&\qquad + \sum_{S \in \mathcal{F} \setminus \mathcal{F}_\mathrm{N}} \int_S [v_h]_S \cdot (\sigma_S - \{\Pi_\Mcal^{k-1} \sigma_\Mcal\}_S \cdot \nu_S) \d{s} + s_h(u_h;v_h).
	\end{align}
	An explicit calculation with
	the definitions of $\sigma_\mathcal{F}$ from \eqref{def:dual-variable} and the stabilization $s_h$ from \eqref{def:stabilization-ldg} proves that the final two terms on the right-hand side cancel. Therefore, $0 = \int_\Omega (\partial_u \psi_h(\bullet, u_h) - \div_h y) v_h \d{x}$ holds for any $v_h \in V_h$ from \eqref{eq:proof-mldg-equivalence}. This yields
	\begin{align}\label{eq:div_h_sigma_h}
		\div_h y = \Pi_\Mcal^k \partial_u \psi_h(\bullet,u_h) = \partial_u \psi_h(\bullet,u_h)
	\end{align}
	under the assumption (B5).
	Since $-sr' + s/(r-1) = -r$, we obtain $\gamma_h(y) = s_h(u_h)$.
	This, the identities $\sigma_\Mcal \cdot \nabla_h u_h = W(\nabla_h u_h) + W^*(\sigma_\Mcal)$ and $\div_h y\,u_h = \psi_h(\bullet,u_h) + \psi_h^*(\bullet,\div_h y)$ from $\nabla_h u_h \in \partial W^*(\sigma_\Mcal)$ and from $u_h \in \partial_u \psi_h^*(\bullet,\div_h y)$ a.e.~in $\Omega$, and \eqref{eq:dELE} with the choice $v_h = u_h$ show
	\begin{align*}
		0 = \int_\Omega (W(\nabla_h u_h) + W^*(\sigma_\Mcal) + \psi_h(x,u_h) + \psi_h^*(x,\div_h y)) \d{x} + \frac{s_h(u_h)}{r} + \frac{\gamma_h(y)}{r'}&.
	\end{align*}
	Rearranging the terms on the right-hand side concludes the proof. 
\end{proof}

\end{document}
